% Generated draft: requires mathematical and source review.
\documentclass[11pt]{article}
\usepackage[margin=25mm]{geometry}
\usepackage{fontspec}
\setmainfont{Latin Modern Roman}
\usepackage{amsmath,amssymb,mathtools,unicode-math}
\setmathfont{Latin Modern Math}
\usepackage{enumitem,needspace,xurl,hyperref}
\setlength{\parindent}{0pt}
\setlength{\parskip}{5pt}
\setlength{\emergencystretch}{3em}
\newcommand{\N}{\mathbb{N}}
\newcommand{\Z}{\mathbb{Z}}
\newcommand{\Q}{\mathbb{Q}}
\newcommand{\R}{\mathbb{R}}
\newcommand{\C}{\mathbb{C}}
\newcommand{\F}{\mathbb{F}}
\newcommand{\A}{\mathbb{A}}
\newcommand{\PP}{\mathbb{P}}
\newcommand{\E}{\mathbb{E}}
\newcommand{\eps}{\varepsilon}
\newcommand{\dd}{\,\mathrm{d}}
\begin{document}
\section*{UnsolvedMath ID: 30002094}
\textbf{Source problem code:} catalogue entry
\subsection{Definitions and mathematical statement}
For specified signal vector spaces \(X\) and \(Y\) over a common scalar field \(K\), let \(S_\tau^X:X\to X\) and \(S_\tau^Y:Y\to Y\) denote compatible time translations for each permitted translation parameter \(\tau\). A map \(T:X\to Y\) is linear and time-invariant if
\[
T(af+bg)=aT(f)+bT(g)
\quad\text{for every }a,b\in K\text{ and }f,g\in X,
\]
and
\[
T S_\tau^X=S_\tau^Y T
\quad\text{for every permitted }\tau.
\]
This gives only the algebraic meaning of linear time-invariance; the supplied record does not specify \(X\), \(Y\), \(K\), or the permitted translations, and does not state whether boundedness is required.

\textbf{Source question.} Characterize the linear time-invariant systems that can be stably approximated by approximation processes using oversampling.

\textbf{Unresolved formulation.} The supplied record does not define the signal class, measurements, approximation processes, sampling threshold relative to which sampling is excessive, or the meaning of stability. It also does not specify a convergence topology, the variables in which convergence must hold uniformly, or the dependence of any stability bound on the input, sampling rate, or approximation index. Consequently, the question above is a characterization request whose precise mathematical predicate remains unspecified. The displayed algebraic definitions do not resolve these omissions or select a particular sampling model.
\subsection{Short English statement}
The question asks which linear time-invariant systems admit stable approximation using oversampling. A precise characterization requires the intended signal spaces, sampling model, approximation procedures, and definitions of stability and convergence. These are absent from the supplied record.
\subsection{Sources}
\begin{itemize}
\item Supplied local UnsolvedMath registry record, permanent ID \(30002094\), title \emph{Stable Oversampling Approximations of Linear Systems}, source problem code \texttt{catalogue entry}, in the Oberwolfach Reports Open Problems collection. The record supplies proposed year \(2012\) and names no proposer.
\item Supplied report citation: \emph{Applied Harmonic Analysis and Sparse Approximation (2013)}, Oberwolfach Reports, DOI \texttt{10.4171/owr/2012/29}, pp.~1759--1843. Source URL and literature-source label ``original OWR report'': \url{https://doi.org/10.4171/owr/2012/29}. The supplied background describes Oberwolfach Reports as documenting mathematical workshops and their open-problem discussions.
\item Supplied literature source: \emph{Signal and system approximation from general measurements}, \url{https://arxiv.org/abs/1402.1092}. No authors or more detailed bibliographic information are supplied in the record.
\item Supplied literature assessment, checked \texttt{2026-08-21}: the record labels the problem open and reports that later work provides stable approximation under broader measurement models and selected sampling schemes, while a general characterization remains unresolved. This is supplied dated curation and literature triage, not an independently verified assessment. Its time-sensitive status and cited claims require independent verification.
\end{itemize}
\end{document}
